% ---------------------------------------------------------------------
% Section 5
% ---------------------------------------------------------------------

\begin{cnread}
静止电荷只有标量势，运动电荷才出现矢势与磁场。本节把「运动」写成 \AEG{}
里的\textbf{方向角} $\theta$：静止时 $\theta=0$（纯加法步，对应只有时间分量
的流），运动时标度与平移同时演化（对应倾斜的世界线）。于是上一节那条
「切触联络 1-形式在轨迹上的配对」$\alpha_1(\dot\gamma_2)$ 立刻给出
$\dfrac{q_1q_2}{4\pi r}\Bigl(1-\dfrac{\mathbf v_1\!\cdot\!\mathbf v_2}{c^2}\Bigr)$：
第二项就是安培力，并且解释了为什么平行电流相吸、为什么 $v\to c$ 时
磁吸引精确抵消静电斥力。本节的符号来源是 Paper~0 定理 3.9；把 $\theta$
认作快度、把 $v$ 认作 $c\tanh\eta$，是本节的第五条字典假设。
\end{cnread}

\subsection{Direction angle and inclined histories}

Paper~0, Theorem~3.9, together with equations (21) and (22) of that paper ---
the angular generator $\Omega_\theta=\mu\cos\theta\,E+\lambda\sin\theta\,H$ and
the flow $\mathrm dz/\mathrm ds=\mu\cos\theta+\lambda z\sin\theta$ --- states
that for a horizontal curve whose tangent is
$\cos\theta\,D_u+\sin\theta\,D_v$ one has
\begin{equation}\label{eq:theta-flow}
  \frac{\mathrm du}{\mathrm ds}=\cos\theta,
  \qquad
  \frac{\mathrm dv}{\mathrm ds}=\sin\theta,
  \qquad
  \frac{\mathrm da}{\mathrm ds}
  =\mu\cos\theta+\lambda a\sin\theta .
\end{equation}
Two readings are distinguished by $\theta$.

\begin{enumerate}[label=(\roman*)]
  \item \emph{At rest}, $\theta=0$: only the additive direction is active,
  $\mathrm dv=0$.  The history is a purely horizontal segment in $\ACS$, the
  assignment evolves by translation, and the current has only the ``time''
  component.  This is the configuration of Section~\ref{sec:coulomb}.
  \item \emph{In motion}, $\theta\ne0$: the same history simultaneously advances
  its scale and its value, $\mathrm dv\ne0$, and in $\ACS$ it is an inclined
  segment whose slope $\mathrm dM/\mathrm d\tau$ measures the motion.  The
  stationary points of the assignment therefore also drift: a moving charge
  changes the scale at which its field is read.
\end{enumerate}

\subsection{The indefinite signature of the frame algebra}

The velocity dictionary is Hypothesis~\ref{hyp:velocity}.  The coefficient
$1-\mathbf v_1\!\cdot\!\mathbf v_2/c^{2}$ that appears below is
not put in by hand: it is what a one-form pairing produces when the vector it is
paired with has an indefinite (Lorentzian) inner product.  Paper~0,
Remark~9.7, is careful to say that the curvature two-form $\omega$ is intrinsic
while its scalar value depends on the chosen horizontal metric; the choice
between a Euclidean and a Lorentzian horizontal metric is therefore exactly the
choice between electrostatics and electrodynamics.

\begin{lemma}[The frame algebra is indefinite]
\label{lem:killing}
In the Lie algebra $\mathfrak{sl}_2(\R)$ with
$X=\begin{pmatrix}m&a_+\\ a_-&-m\end{pmatrix}$ one has
$\det X=-(m^{2}+a_+a_-)$.  Consequently the quadratic form
$(m,a_+,a_-)\mapsto\det X$ has signature $(2,1)$: the direction of the scale
rate $m$ enters with the opposite sign to the mixed translation entries.
\end{lemma}

\begin{proof}
$\det X=-m^{2}-a_+a_-$ is a direct computation.  The form $-m^{2}-a_+a_-$ is
negative definite on the line $a_\pm=0$ and of split sign on the plane
$m=0$, whence signature $(1,2)$ up to an overall sign, i.e. $(2,1)$ after
overall negation.
\end{proof}

\begin{remark}
Lemma~\ref{lem:killing} is the algebraic shadow of the statement that the time
direction and the space directions of a moving source cannot be treated on the
same footing.  It is recorded here because it explains why the coefficient in
\eqref{eq:amperes-coefficient} is a difference rather than a sum.  No
identification of $m$ with a physical time direction is made; that identification
is Hypothesis~\ref{hyp:velocity}.
\end{remark}

\subsection{The contact pairing for two moving charges}

Let charge 1, with charge $q_1$ and velocity $\mathbf v_1$, excite the gauge
potential read as the contact one-form
\begin{equation}\label{eq:gauge-potential}
  \alpha_1=\phi_1\,\mathrm dt-\mathbf A_1\!\cdot\!\mathrm d\mathbf x
  =\frac{q_1}{4\pi r}\Bigl(\mathrm dt-\mathbf v_1\!\cdot\!\mathrm d\mathbf x\Bigr).
\end{equation}
In \AEG{} terms, $\mathrm dt$ is the pure additive frame direction and
$\mathbf v_1\!\cdot\!\mathrm d\mathbf x$ is the synchronous drift of the scale
coordinate caused by the motion, i.e. the $\beta_1$ multiple of the
multiplicative direction of Hypothesis~\ref{hyp:velocity}.  Let charge 2 follow
its own inclined history with tangent $\dot\gamma_2=(1,\mathbf v_2)$.

\begin{proposition}[Amp\`ere coefficient from a one-form pairing]
\label{prop:ampere}
Under Hypotheses~\ref{hyp:charge}--\ref{hyp:velocity}, the interaction
Lagrangian density of the two charges is the evaluation of the contact form on
the second history's tangent,
\begin{equation}\label{eq:amperes-lagrangian}
  L_{\mathrm{int}}=q_2\,\alpha_1(\dot\gamma_2)
  =\frac{q_1q_2}{4\pi r}\Bigl(1-\frac{\mathbf v_1\!\cdot\!\mathbf v_2}{c^{2}}\Bigr),
\end{equation}
so that the effective potential and force are
\begin{equation}\label{eq:amperes-coefficient}
  V_{\mathrm{eff}}(r)=\frac{q_1q_2}{4\pi r}
   \Bigl(1-\frac{\mathbf v_1\!\cdot\!\mathbf v_2}{c^{2}}\Bigr),
  \qquad
  F(r)=-\frac{\partial V_{\mathrm{eff}}}{\partial r}
  =\frac{q_1q_2}{4\pi r^{2}}
   \Bigl(1-\frac{\mathbf v_1\!\cdot\!\mathbf v_2}{c^{2}}\Bigr).
\end{equation}
\end{proposition}

\begin{proof}
Insert $\dot\gamma_2=(1,\mathbf v_2)$ into \eqref{eq:gauge-potential} to get
$\alpha_1(\dot\gamma_2)=\frac{q_1}{4\pi r}\bigl(1-\mathbf v_1\!\cdot\!\mathbf
v_2\bigr)$ (with $c$ restored in the second term), and multiply by the second
charge.  This is the interaction Lagrangian density; for a static pair it
coincides with the potential energy of \eqref{eq:pairing-energy} at
$\mathbf v_1=\mathbf v_2=0$, so no sign change is needed in passing to
$V_{\mathrm{eff}}$.  Differentiating in $r$ gives
\eqref{eq:amperes-coefficient}.
\end{proof}

The coefficient splits into the static Coulomb part of
Section~\ref{sec:coulomb} and a purely motional part:
\begin{equation}\label{eq:ampere-split}
  F(r)=\underbrace{\frac{q_1q_2}{4\pi r^{2}}}_{\text{static}}
       \;-\;\underbrace{\frac{q_1q_2\,(\mathbf v_1\!\cdot\!\mathbf v_2)}{4\pi c^{2}r^{2}}}_{\text{motional}} .
\end{equation}

\begin{corollary}[Parallel currents attract the electrons]
\label{cor:parallel-currents}
For two electrons, $q_1q_2=+e^{2}$, moving in parallel so that
$\mathbf v_1\!\cdot\!\mathbf v_2=v^{2}>0$, the motional term is negative and
reduces the repulsion.  As $v\to c$ it cancels the static term exactly and the
total force tends to zero.
\end{corollary}

\begin{proof}
Immediate from \eqref{eq:ampere-split}: the bracket becomes
$1-v^{2}/c^{2}\to0$.
\end{proof}

The sign structure of \eqref{eq:ampere-split} is the familiar one: like charges
repel electrostatically and attract magnetically when their currents are
parallel, and \eqref{eq:amperes-coefficient} is the sum.  What the \AEG{} reading
adds is not the algebra --- which is the standard tree-level Breit--Darwin
coefficient --- but the fact that the coefficient arises as a single pairing of
one contact one-form with one inclined tangent, rather than as a sum of two
separately postulated field contributions.

\begin{warning}[Range of validity]
\label{warn:darwin}
Equation~\eqref{eq:amperes-coefficient} is the quasi-static coefficient of order
$v^{2}/c^{2}$; in electrodynamics it is the Darwin term of the two-body
Lagrangian \cite{Darwin1920,Jackson1999}.  It is not the exact force between two
uniformly moving charges: the complete expression has further velocity-dependent
terms of order $v^{4}/c^{4}$, retardation corrections, and spin--orbit
contributions.  Nothing in this section supplies those terms, and no
all-orders statement is claimed.  In the reading of
Hypothesis~\ref{hyp:velocity} the truncation corresponds to the linearised
contact flow \eqref{eq:theta-flow}, which is exact for constant $\theta$ but
does not by itself generate the full Lorentz-covariant field.
\end{warning}

\subsection{Uniform motion in ripple geometry: compression without radiation}

Paper~0, Section~4, reads the same word contravariantly as a family of pulled-back
horocycles.  For a prefix $P_k$, the ripple pencil is the family
$\mathcal R(P_k)=\{P_k^{-1}(\mathcal H_t)\}$ of circles tangent at the pole
$p_k=P_k^{-1}(\infty)$ (Theorem~4.2).  The two cases now read as follows.

\begin{enumerate}[label=(\roman*)]
  \item \emph{Charge at rest.}  The pole is fixed, and the wavefronts of
  successive times are concentric or commonly tangent circles.  The resulting
  field has the $1/r^{2}$ shape of a static configuration.
  \item \emph{Charge in uniform motion.}  A Lorentz boost is a hyperbolic
  M\"obius transformation, and the celestial sphere's Lorentz group is
  $\PSL_2(\CC)$, whose two-dimensional section is $\PSL_2(\R)$ --- the \AEG{}
  M\"obius group itself.  Uniform motion therefore acts on the ripple pencil by
  a hyperbolic element $D_{e^{\eta}}\in H$, compressing the wavefronts ahead and
  stretching them behind, while leaving the pencil's common tangency point fixed:
  the two bivaluation data $(q^+,q^-)$ are unchanged, so the history stays inside
  a single stabilizer orbit (Paper~IV, Theorem~3.5 and
  Corollary~4.4, together with Proposition~4.6 for the normalizer).
\end{enumerate}

The second case is the geometric form of the statement that a uniformly moving
charge carries its near field with it and does not shed radiation: the pencil is
compressed, not disrupted.  Section~\ref{sec:radiation} proves the corresponding
algebraic statement --- that a closed history with constant kernel coordinate has
holonomy exactly $1$ --- and shows that the statement fails as soon as the kernel
coordinate moves.

\begin{remark}[The $v\to c$ limit in the $\E_0$ picture]
\label{rem:v-to-c}
Corollary~\ref{cor:parallel-currents} can be read in $\E_0$ as follows.  As the
common tilt of two histories approaches the critical rate, their tangent
directions approach a common asymptotic slope, i.e. the two curves approach a
common ideal point.  In the picture of Paper~0, Theorem~3.9, the relative rate of
shear between them decreases like $1-\beta^{2}$, so the relative transverse
disruption of the two pencil families vanishes.  This is a heuristic reading of
the coefficient, not a derivation of it: the derivation is
Proposition~\ref{prop:ampere}, and the heuristic is included because it is what
makes the coefficient visible inside $\E_0$.
\end{remark}
